\section{Reaching both original growth targets with retained diffusion}
\label{ave:section}

We now join the exact evolving-parent growth calculation to the direct
second return theorem. Both stopping events are the original temperature
thresholds. This changes the first growth time from the earlier
fixed-inviscid-time variant in Section~\ref{sec:cv-diffusive-parent}.
The construction below specifies the frequencies and a nonempty interval
of positive physical diffusivities, and proves that their actual entry
states satisfy every inequality in Theorem~\ref{sr:theorem}.
The spatial fields and their full forces are supplied by the compact
construction following that theorem. In particular the amplitude damping
on the released linear profile has its explicitly stated core force.

\subsection{A finite choice within the original frequency schedule}

Retain the original physical data
\[
 A_0\geq1,\qquad\nu=\sqrt{A_0},\qquad
 \mu=\lambda_0=\lceil\sqrt{A_0}/2\rceil,
\]
and all fixed source profiles, including
$s_*,s_0,C_\ell,\eta,\eta_2,\Lambda_1$.
Thus $0<s_*\leq1/8$, $C_\ell\geq4$, $H=\|\eta_2\|_\infty\geq1$,
$\Lambda_1\geq6H$, and $3\Lambda_1H\sin s_*\leq1/4$.
The first section's $\lambda_1$ denoted a source frequency; in this section
we write that same quantity as $\widehat\lambda_1$ and retain the
physical frequencies separately as $K_j=\mu\widehat\lambda_j$.
Keep the source's $Q_*\geq200$, $Q_j=Q_*+j$, its specified finite
constants $K(k)$, and the exact choices
\begin{align}
 \widehat\lambda_2&=\widehat\lambda_1^{Q_2},&
 k_j&=\max\{k\in\mathbb N_0:120k\leq Q_j,
                   \ K(k)\leq\log\widehat\lambda_j\},\nonumber\\
 L_j&=(k_j+41/8)\log\widehat\lambda_j,&
 S_j&=\frac{A_0}{\mu}\widehat\lambda_j^{-k_j-6},\nonumber\\
 P_j^*&=\frac{A_0}{\mu}\widehat\lambda_j^{-7/8}=S_je^{L_j}.
 \label{ave:source_choices}
\end{align}
Here $K(k)$ is the source's prescribed parameter sequence, not the
physical frequency $K_j$. The maximum is a nonempty finite set once
$\log\widehat\lambda_1\geq K(0)$: $k=0$ belongs to it and
$k\leq\lfloor Q_j/120\rfloor$. No smooth dependence of $k_j$ on the
frequencies is required.

The following fixed constants give one explicit sufficient frequency
choice; their large numerical size is immaterial to finite existence:
\begin{equation}\label{ave:Ychoice}
\begin{gathered}
 C_A=e^{2+\Lambda_1^{-1}},\qquad
 C_\sigma=\sqrt{1+C_A},\qquad C_N=1+3C_\sigma/32,\\
 C_L=(\lfloor Q_2/120\rfloor+41/8)Q_2,\\
 Y_0=1+\max\left\{
 \begin{gathered}
 64,K(0),Q_2,72\cdot16^3HC_L^2/\sqrt e,\quad8\log(4/e),\\
 (16/47)\log(4C_\ell C_N/s_0),\quad
 (16/31)\log(2C_N/s_0)
 \end{gathered}\right\}.
\end{gathered}
\end{equation}
Choose any $\widehat\lambda_1\geq e^{Y_0}$ and then use
\eqref{ave:source_choices} literally. Write $L=L_2$,
$\Gamma_1=\nu\sin s_*$ and $C_{\rm post}=2+\Lambda_1^{-1}$.
Define the explicit positive number
\begin{equation}\label{ave:dmax}
 d_{\max}=\min\left\{
 \frac{\Gamma_1}{2K_1^2},
 \frac{\Gamma_1\log2}{K_1^2C_{\rm post}},
 \frac{\nu\log(3/2)}{16K_1^2},
 \frac{L\nu\log3}{12800K_2^2}\right\}.
\end{equation}
Every entry in this minimum uses only the fixed original data and
the chosen frequencies. In particular it does not require solving
a circular inequality involving a yet unknown diffusive growth scale.
All subsequent statements hold for each $0<d\leq d_{\max}$, with
physical viscosity and thermal diffusivity both equal to $d$.

\subsection{First actual threshold and the changed first endpoint}

Put $q=dK_1^2$ and $t_1^{\rm in}=1/\nu$.
Growth and transition still have $\alpha=0$. The signed first pair is
exactly
\begin{equation}\label{ave:first_growth}
 \Theta_1(t)=-S_1e^{(\Gamma_1-q)(t-t_1^{\rm in})},\qquad
 \Omega_1(t)=\frac{K_1}{\nu}\Theta_1(t).
\end{equation}
Indeed its temperature coupling is $a_1=A_0\sin s_*/K_1$,
its vorticity coupling is $b_1=K_1\sin s_*$, and
$a_1K_1/\nu=\Gamma_1$, $b_1\nu/K_1=\Gamma_1$.
Differentiation verifies both damped equations and both original
seed values. The linear integral equation proves uniqueness.
Since $q\leq\Gamma_1/2$, the original first threshold is attained
uniquely and transversely at
\begin{equation}\label{ave:first_times}
 t_1^a=t_1^{\rm in}+\frac{L_1}{\Gamma_1-q},\qquad
 t_1^1=t_1^a+\Gamma_1^{-1},\qquad
 t_1^b=t_1^1+\Gamma_1^{-1}(1+\Lambda_1^{-1}).
\end{equation}
The first activation $\eta(\Gamma_1(t-t_1^{\rm in}))$ is already one
at $t_1^a$, since $L_1\geq1$.

Starting at $t_1^a$, take the original first transition and steering
control with their displayed durations, translated to these actual
times. The inverse map
\[
 (\Theta_1,\Omega_1)(t)
   =e^{-q(t-t_1^a)}(\Theta_{1,E},\Omega_{1,E})(t)
\]
relates their damped solution to the inviscid controlled solution
with data
\[
 (\Theta_{1,E},\Omega_{1,E})(t_1^a)
       =(-P_1^*,-K_1P_1^*/\nu).
\]
Differentiation gives the two damped equations, and multiplication
by $e^{q(t-t_1^a)}$ gives the inverse equations and identical data.
Both factors are positive. Thus the quotient, its entire trial
family, and its unique first return root proved in
Section~\ref{sec:cv-first-stage} are unchanged on these translated
intervals. At the selected endpoint
\begin{equation}\label{ave:first_endpoint}
 \alpha=s_*,\quad\zeta_1=e_2,\quad\Omega_1=0,\quad
 A_d:=-K_1\Theta_1(t_1^b)
       =A_Ee^{-qC_{\rm post}/\Gamma_1},
\end{equation}
where the proved inviscid endpoint satisfies
\begin{equation}\label{ave:AE_bounds}
 2eA_0\widehat\lambda_1^{1/8}\leq A_E
          \leq C_A A_0\widehat\lambda_1^{1/8}.
\end{equation}
This expression does not damp the inviscid growth duration: that
duration has been replaced by the actual threshold time in
\eqref{ave:first_times}. The first three diffusion exponents are
exactly $qL_1/(\Gamma_1-q)$ in growth, $q/\Gamma_1$ in transition,
and $q(1+\Lambda_1^{-1})/\Gamma_1$ in steering. The physical gains
are respectively $L_1$, $1-q/\Gamma_1$, and the original selected
steering gain minus the last exponent.

\subsection{Actual second scale and all fixed design inequalities}

The true inherited gradient at $t_2^{\rm in}=t_1^b$ is
$G_d=(A_0\sin s_*,-A_0\cos s_*-A_d)$. Evaluate the source scale
at this changed endpoint, as required by its original rule:
\begin{equation}\label{ave:actual_scale}
 \sigma^2=|G_d|
   =\sqrt{A_0^2+A_d^2+2A_0A_d\cos s_*},\quad
 s=s_2=\frac{L\nu}{\sigma},\quad a=\sin s,\quad
 \Gamma=\Gamma_2=\sigma a.
\end{equation}
Since $qC_{\rm post}/\Gamma_1\leq\log2$, equations
\eqref{ave:AE_bounds}--\eqref{ave:actual_scale} imply
\begin{equation}\label{ave:scale_bounds}
 eA_0\widehat\lambda_1^{1/8}\leq A_d\leq
 C_A A_0\widehat\lambda_1^{1/8},\qquad
 \sqrt e\,\nu\widehat\lambda_1^{1/16}\leq\sigma
       \leq C_\sigma\nu\widehat\lambda_1^{1/16}.
\end{equation}
For the lower scale bound use $\sigma^2\geq A_d$, since the square
of its left side differs from $A_d^2$ by a positive quantity.
For the upper bound use the triangle inequality
$\sigma^2\leq A_0+A_d\leq(1+C_A)A_0\widehat\lambda_1^{1/8}$.
The choice $Y_0>8\log(4/e)$ also gives $A_d>4A_0$.

Let $y=\log\widehat\lambda_1\geq Y_0$.
Then $64<L\leq C_Ly$. The exponential series with nonnegative
terms proves $e^{y/16}\geq y^3/(6\cdot16^3)$; by
\eqref{ave:Ychoice} its right side is strictly greater than
$12HC_L^2y^2/\sqrt e$. Consequently
\begin{equation}\label{ave:angle_bounds}
 0<s\leq\frac{L}{\sqrt e\,e^{y/16}}
      <\frac1{12HL}\leq\frac18,
 \qquad 0<a\leq\frac18,
 \qquad 3LHa<\frac14.
\end{equation}
For $0\leq s\leq1/8$, integrating $\cos r\geq1-r^2/2\geq1/2$
on $[0,s]$ gives $s/2\leq\sin s\leq s$. Thus
\begin{equation}\label{ave:Gamma_bounds}
 \frac{L\nu}{2}\leq\Gamma\leq L\nu.
\end{equation}
The fixed base coefficients in the second-return coordinates obey
\begin{equation}\label{ave:C_bounds}
 0<C_0=\frac{A_0\cos s_*}{\sigma^2}<\frac14,
 \qquad 0<C_1=\frac{A_0\sin s_*}{\sigma^2a}
 \leq\frac{2\nu\sin s_*}{\sigma L}<\frac1{32}.
\end{equation}
The last inequality uses $\sigma\geq\nu$, $L\geq64$ and
$\sin s_*<1$; it has no missing comparison constant.

The same frequency choice proves the source support nesting for the
entire second return and its finite hold. Set
$N=1+(3-\cos s)/a$. The exact matrix reconstruction in
Theorem~\ref{sr:theorem} gives both $\|M_2\|$ and $\|M_2^{-1}\|$
at most $N$, by the triangle inequality for its displayed shear.
Equations~\eqref{ave:scale_bounds} and $a\geq s/2$ imply
$1/a\leq 2C_\sigma\widehat\lambda_1^{1/16}/L$, whence
$N\leq C_N\widehat\lambda_1^{1/16}$.
The final two logarithmic bounds in \eqref{ave:Ychoice} give, strictly,
\begin{equation}\label{ave:nesting}
 N\widehat\lambda_1^{-3}<\frac{s_0}{2C_\ell},
 \qquad N\widehat\lambda_1^{-2}<s_0.
\end{equation}
Indeed their respective left sides are at most
$C_Ne^{-47y/16}$ and $C_Ne^{-31y/16}$, and our choices bound these
by $s_0/(4C_\ell)$ and $s_0/2$.
For $\ell_1=s_0/(C_\ell\mu)$ and
$\ell_2=\widehat\lambda_1^{-3}/\mu$, this places the entire second
support inside the first envelope's constant region and its exact
linear profile region: $|x|\leq N\ell_2<\ell_1/2$ and
$|K_1\zeta_1\cdot x|\leq K_1N\ell_2<s_0$.
Every original factor of $\mu$ is present in these identities.

\subsection{Second growth with the actual diffusing parent}

Write $r=t-t_2^{\rm in}$. The second growth and transition have
$\alpha=s_*$, $\Omega_1=0$, $M_2=I$, $\zeta_2=e(s)$ and
$A_1(r)=A_de^{-qr}$. These claims follow from the closed subsystem:
$\zeta_{1,1}=0$ makes $\dot\Omega_1=-q\Omega_1=0$;
both controls $F_1^{\rm ctl}$ and $F_2^{\rm ctl}$ are then zero,
and the phase matrix equation is $\dot M_2=0$. The first temperature
equation is $\dot A_1=-qA_1$.
The second pair has the exact coefficients
\begin{equation}\label{ave:growth_coefficients}
 a_2(r)=\frac{A_de^{-qr}\sin s+A_0\sin(s_*+s)}{K_2},
 \qquad b_2=K_2\sin s,
 \qquad d_2^{\rm phys}=dK_2^2,
\end{equation}
and the exact original entry data
$\Theta_2(0)=-S_2$, $\Omega_2(0)=\sqrt{b_2/a_2(0)}\Theta_2(0)$.
This is precisely the system solved by the convergent fundamental
functions in Section~\ref{sec:cv-diffusive-parent}; the substitution
here uses its displayed changed coefficients, not frozen values.

Put $a_i=a_2(0)$, $u_i=\sqrt{b_2/a_i}$,
$\gamma_i=\sqrt{a_i b_2}$ and
\[
 n_i=\frac{\gamma_i^2}{\Gamma^2}
 =\frac{A_d+A_0\cos s_*}{\sigma^2}+C_1\cos s.
\]
One has
\begin{equation}\label{ave:ni_bounds}
 \frac45<n_i<\frac{33}{32}<2.
\end{equation}
The lower bound follows from
$A_d/\sigma^2\geq A_d/(A_d+A_0)>4/5$.
For the upper bound, square
$\sigma^2=\sqrt{(A_d+A_0\cos s_*)^2+A_0^2\sin^2s_*}$
to get $(A_d+A_0\cos s_*)/\sigma^2\leq1$, then use
\eqref{ave:C_bounds}. In particular
$\gamma_i\geq\Gamma/\sqrt2\geq L\nu/(2\sqrt2)$.

On the explicitly fixed horizon $0\leq r\leq16/\nu$, the third
bound in \eqref{ave:dmax} gives $qr\leq\log(3/2)$. Since both
terms in $a_2$ are positive, it follows that
\begin{equation}\label{ave:coeff_interval}
 \frac12a_i\leq a_2(r)\leq a_i.
\end{equation}
To prove all signs and ratio estimates without assuming them, set
$X=-e^{d_2^{\rm phys}r}\Theta_2$ and
$Y=-e^{d_2^{\rm phys}r}\Omega_2$. They solve
$X'=a_2Y$, $Y'=b_2X$ with positive initial values.
Before any first zero, both derivatives are positive, excluding that
zero by continuity. The linear integral equation, whose coefficient
matrix is bounded on this horizon, gives existence there by its
uniformly convergent iterated-integral series.
Therefore $u=Y/X=\Omega_2/\Theta_2$ exists and satisfies
$u'=b_2-a_2u^2$. The constant functions $u_i$ and $\sqrt2u_i$
are respectively subsolution and supersolution on this horizon,
because $b_2-a_2u_i^2\geq0$ and
$b_2-2a_2u_i^2\leq0$. Subtracting either comparison equation
from the equation for $u$ gives the linear difference coefficient
$-a_2(u+u_{\rm comparison})$; its integrating factor preserves
the stated ordering. Hence
\begin{equation}\label{ave:ratio_growth}
 u_i\leq u\leq\sqrt2u_i,\qquad
 \frac1{\sqrt{n_i}}\leq
 v:=\frac{\sigma\Omega_2}{K_2\Theta_2}
 \leq\sqrt{\frac2{n_i}},\qquad\frac12<v<2.
\end{equation}
Here $\sigma u_i/K_2=\Gamma/\gamma_i=1/\sqrt{n_i}$;
the equality follows by substituting the original $a_i,b_2,K_2$.

For $P_2=-\Theta_2$, the exact physical logarithmic derivative is
\[
 (\log P_2)'=a_2u-d_2^{\rm phys}.
\]
The last bound in \eqref{ave:dmax} gives
$d_2^{\rm phys}\leq L\nu\log3/12800<L\nu/16
\leq\gamma_i/4$. Thus throughout the fixed horizon
\begin{equation}\label{ave:growth_rate}
 \frac{\gamma_i}{4}\leq(\log P_2)'
      \leq\sqrt2\gamma_i\leq2\Gamma.
\end{equation}
Indeed $a_2u\geq a_iu_i/2=\gamma_i/2$ and
$a_2u\leq\sqrt2a_iu_i=\sqrt2\gamma_i$.
Consequently the original second target $P_2^*=S_2e^L$ is reached
at exactly one $r=r_a$, with nonzero derivative, and
\begin{equation}\label{ave:second_times}
 \frac{L}{2\Gamma}\leq r_a\leq\frac{4L}{\gamma_i}
       \leq\frac{8\sqrt2}{\nu},\qquad
 r_1=r_a+\Gamma^{-1}<\frac{16}{\nu},
 \quad t_2^a=t_2^{\rm in}+r_a,
 \quad t_2^1=t_2^{\rm in}+r_1.
\end{equation}
For existence of the crossing, the upper candidate time is inside
the fixed horizon; integrating the lower rate to that time gives
at least $L$ logarithmic gain. Strict monotonicity and the intermediate
value theorem give existence and uniqueness. The transition length
is at most $2/(L\nu)\leq1/(32\nu)$, so the entire transition is
also inside that already justified horizon. The lower bound on $r_a$
ensures the second activation $\eta(\Gamma r)$ equals one before
the growth threshold, since $L\geq64$.

\subsection{Closing every entry bound and the complete finite statement}

At the actual steering entry $t_2^1$, take
$P_{1\rm in}=A_de^{-qr_1}/K_1$ and
$P_{2\rm in}=P_2(r_1)$. Equations~\eqref{ave:ratio_growth} and
\eqref{ave:second_times} give the initial quotient interval in
Theorem~\ref{sr:theorem}. Also
\begin{equation}\label{ave:Bentry}
 \frac{8}{15}<B_{\rm in}
 =\frac{A_de^{-qr_1}}{\sigma^2}\leq1,
 \qquad
 0<\delta_2=\frac{dK_2^2}{\Gamma}\leq\frac{\log3}{6400},
 \qquad 0<\delta_1=\frac{dK_1^2}{\Gamma}
       \leq\delta_2<\frac1{16}.
\end{equation}
The lower parent bound is $(4/5)(2/3)$; the upper is
$A_d\leq\sigma^2$. The diffusion bounds follow from
\eqref{ave:dmax}, \eqref{ave:Gamma_bounds} and $K_1<K_2$.
Equations~\eqref{ave:angle_bounds}, \eqref{ave:C_bounds},
\eqref{ave:ratio_growth}, and \eqref{ave:Bentry} prove every
entry/design hypothesis of the direct return theorem from the original
seeds and the explicit choices \eqref{ave:Ychoice}--\eqref{ave:dmax}.

\begin{theorem}\label{ave:complete_entry}
For each fixed original source datum and profiles above, choose
$\widehat\lambda_1$ by \eqref{ave:Ychoice}, retain the source recursion
\eqref{ave:source_choices}, and choose any $0<d\leq d_{\max}$ in
\eqref{ave:dmax}. The complete modified two-layer coefficient system
exists through both original growth targets, both prescribed
transitions, both selected vorticity returns, and a second holding
interval of physical length $1/(32\Gamma_2)$.
Its times, amplitudes, original coordinates, controls, and spatial
support inequalities are the exact formulas above and the inverse
formulas of Theorem~\ref{sr:theorem}. At the second return the new
vorticity amplitude and horizontal phase component are zero, the
first parent vorticity is strictly positive, and
\begin{equation}\label{ave:total_gain}
 \log\frac{-\Theta_2(t)}{S_2}
 \geq L_2+\frac{\gamma_i}{4\Gamma_2}+\frac{\log3}{320}
 \quad\text{for }t_2^b\leq t\leq t_2^b+(32\Gamma_2)^{-1}.
\end{equation}
The inherited parent is evolved throughout; no inviscid estimate is
assumed to persist under diffusion. With the full profile/localization
forces below these coefficients determine compact spatial solutions
from the original initial datum and zero initial velocity.
\end{theorem}
\begin{proof}
The first stage was proved by the exact map at its actual threshold.
The second growth and transition, their signs and all entry bounds
were proved on a common explicit physical horizon. Theorem~\ref{sr:theorem}
therefore applies, constructs its entire trial family before taking
the least root, and proves the return and nonstationary parent hold.
The second growth contributes exactly $L_2$ to logarithmic gain.
The transition contributes at least $\gamma_i/(4\Gamma_2)$ by
\eqref{ave:growth_rate}. The steering and entire hold contribute at
least $\log3/320$ by that theorem and \eqref{ave:Bentry}.
Adding these three actual integrals proves \eqref{ave:total_gain}.
All finite join controls are smooth by flatness of the original step
and interior support of its pulse; the solution of the smooth ODE
therefore matches to every order at those joins. The spatial identities
below prove the stated compact realization with every required force.
\end{proof}

For use in those spatial bounds, throughout both stages and their hold
one may use the following proved amplitude upper bounds:
\begin{align}
 \sup|\Theta_1|&\leq
       \max\{C_AP_1^*,2\sigma^2/K_1\},&
 \sup|\Omega_1|&\leq
       \max\{(K_1/\nu)C_AP_1^*,9\sigma\},\nonumber\\
 \sup|\Theta_2|&\leq e^{18}P_2^*,&
 \sup|\Omega_2|&\leq (3K_2/\sigma)e^{18}P_2^*.
 \label{ave:amplitude_sup}
\end{align}
For layer one, use its first-stage positive multiplier bounds and
the direct second-stage bounds $B\leq2,W<9$.
For layer two, growth never exceeds its threshold; the transition
has dimensionless coefficient $a_2K_2/(\sigma\Gamma)
=a_2b_2/\Gamma^2\leq n_i<2$ and $v<2$, so its logarithmic
gain is at most $4$. The selected steering gain is at most
$12(1+L^{-1})<14$, and holding only decreases temperature.
These observations prove $e^{18}P_2^*$ as well as the vorticity
bound from $0\leq v\leq3$ on the selected steering, $v<2$ before it,
and $v=0$ on holding.

The exact second-stage diffusion exponents before steering are
$dK_2^2r_a$ and $dK_2^2/\Gamma$; the first parent's exponent
over these same intervals is $q(r_a+1/\Gamma)$.
The steering and holding exponents retain the evolving $|\zeta_2|^2$
inside the integrals in Theorem~\ref{sr:theorem}; together these are
the complete two-stage amplitude diffusion costs.

Finally the proved viscosity range is explicitly frequency dependent.
For fixed original data and $Q_2$,
\begin{equation}\label{ave:range_limit}
 0<d_{\max}\leq
 \frac{\nu C_L\log3}{12800\mu^2}
 (\log\widehat\lambda_1)\widehat\lambda_1^{-2Q_2}
 \longrightarrow0\qquad(\widehat\lambda_1\longrightarrow\infty).
\end{equation}
To verify the limit, set $y=\log\widehat\lambda_1$ and use
$e^{2Q_2y}\geq(2Q_2y)^2/2$, giving
$ye^{-2Q_2y}\leq1/(2Q_2^2y)\to0$.
This is the dependence of the range proved here, not an assertion
that every larger diffusivity fails or that a modified infinite
construction is impossible. The actual finite result preserves both
original thresholds and controls. A fixed-positive-diffusivity
iteration would require further estimates or a changed scale design;
it does not follow by iterating a range whose upper bound tends to zero.

\subsection{An exact scaling restriction for this growth rule}

The original datum ties the physical spatial scale to the initial
amplitude: $\mu=\lceil\nu/2\rceil$ with $\nu=\sqrt{A_0}\geq1$.
Consequently
\begin{equation}\label{ave:original_scaling_restriction}
 \frac{\mu^2}{\nu}\geq\frac12,
 \qquad \widehat d=d\frac{\mu^2}{\nu}\geq\frac d2,
\end{equation}
where $\widehat d$ is exactly the transformed diffusion coefficient
in the original map $(X,T)=(\mu x,\nu t)$.
To prove the first inequality, if $1\leq\nu\leq2$ then $\mu=1$
and $\mu^2/\nu=1/\nu\geq1/2$. If $\nu>2$, then
$\mu\geq\nu/2$, whence $\mu^2/\nu\geq\nu/4>1/2$.
Equality occurs at the original datum value $A_0=4$, $\nu=2$, $\mu=1$;
this identifies the equality case and does not change the parameters
of any constructed trajectory. The second inequality follows by
multiplication by $d>0$, using the exact diffusion operator map.

For the first modified growth rule proved here, attaining
$P_1^*=S_1e^{L_1}>S_1$ requires and is equivalent to
$dK_1^2<\Gamma_1$. This follows directly from
\eqref{ave:first_growth}: at equality the temperature amplitude is
constant, and above it the amplitude strictly decreases for all growth
times. Therefore every instance of this same growth rule must satisfy
\begin{equation}\label{ave:necessary_first_growth}
 \widehat d\,\widehat\lambda_1^2<\sin s_*,\qquad
 d\widehat\lambda_1^2<2\sin s_*.
\end{equation}
The first inequality is the identical physical inequality divided by
$\nu$, since $K_1=\mu\widehat\lambda_1$; the second follows from
\eqref{ave:original_scaling_restriction}.
Thus increasing $A_0$ while preserving the original spatial scaling
cannot make $\widehat d$ tend to zero at a fixed positive physical $d$.
This necessary first-growth restriction is a result about the specified
modified amplitude equation and original growth control. Any route
past it must change an identified part of that dynamics or scale
design and verify the corresponding full spatial forces. It is not
an obstruction to all forced viscous flows.
